\documentclass[a4paper,12pt]{article}
\usepackage[T2A]{fontenc}
\usepackage[utf8]{inputenc}
\usepackage[russian]{babel}
\usepackage{amsmath,amssymb,mathtools}
\usepackage{helvet}
\renewcommand{\familydefault}{\sfdefault}
\usepackage{icomma}
\usepackage[a4paper,top=19mm,bottom=20mm,left=20mm,right=20mm]{geometry}
\usepackage{microtype}
\usepackage{tikz}
\usetikzlibrary{calc,arrows.meta,positioning,shapes.geometric}
\usepackage{tabularx,booktabs}
\usepackage{enumitem}
\usepackage{parskip}
\usepackage{fancyhdr}
\usepackage{setspace}
\usepackage{xcolor}
\usepackage{hyperref}

\definecolor{accent}{HTML}{256C70}
\definecolor{accentdark}{HTML}{174C52}
\definecolor{ink}{HTML}{242C30}
\definecolor{muted}{HTML}{68757B}
\definecolor{hair}{HTML}{CDD7D7}
\hypersetup{hidelinks}
\color{ink}
\setlength{\parindent}{0pt}
\setlength{\parskip}{5pt}
\setlist[itemize]{leftmargin=6mm,topsep=3pt,itemsep=2pt}
\setlist[enumerate]{leftmargin=8mm,topsep=3pt,itemsep=5pt}
\onehalfspacing
\pagestyle{fancy}
\fancyhf{}
\renewcommand{\headrulewidth}{0pt}
\fancyfoot[C]{\footnotesize\color{muted}09.10.26\quad\textbullet\quad Марк Гукин\quad\textbullet\quad\thepage}
\newcommand{\chap}[1]{\par\medskip\noindent{\Large\bfseries\color{accentdark}#1}\par\vspace{2pt}{\color{hair}\rule{\linewidth}{0.45pt}}\par\vspace{3pt}}
\newcommand{\subchap}[1]{\par\smallskip{\color{accentdark}\bfseries #1}\par}
\newcommand{\example}[1]{\par\smallskip\textbf{#1}\par}
\newcommand{\key}[1]{\textcolor{accentdark}{\textbf{#1}}}
\newcommand{\mymath}[1]{\[\displaystyle #1\]}

\begin{document}
\thispagestyle{empty}
\begin{center}
\vspace*{25mm}
{\large\color{muted}МАТЕМАТИКА\quad /\quad ПОДГОТОВКА К ЕГЭ}\par
\vspace{11mm}
{\fontsize{32}{38}\selectfont\bfseries\color{accentdark}Чек-лист}\par
\vspace{7mm}
{\fontsize{20}{26}\selectfont\bfseries Показательные и степенные уравнения}\par
\vspace{3mm}
{\fontsize{17}{23}\selectfont Прикладные задачи по формуле}\par
\vspace{11mm}
{\normalsize Свойства степеней \textbullet\; примеры \textbullet\; алгоритм \textbullet\; 10 задач}\par
\vfill
{\large 09.10.26}\par
\vspace{8mm}
{\normalsize Лёвин Артём Александрович}\par
{\normalsize\bfseries эксклюзивно для Марка Гукина}\par
\vspace{13mm}
\end{center}
\begin{tikzpicture}[remember picture,overlay]
\draw[accent!55,line width=1.15pt]
  ([xshift=19mm,yshift=17mm]current page.south west)
  .. controls +(44mm,14mm) and +(-51mm,-7mm) ..
  ([xshift=-19mm,yshift=20mm]current page.south east);
\end{tikzpicture}
\clearpage

\chap{1. Инструменты: степени и корни}
Перед решением приведите составные числа, дроби и корни к степеням с удобным основанием. Это общий приём для всех трёх блоков занятия.
\mymath{8=2^3,\quad 16=2^4,\quad 125=5^3,\quad
\frac12=2^{-1},\quad \sqrt[5]{16}=2^{4/5}.}
Для $a>0$ и действительных показателей $u,v$:
\begin{align*}
a^u\cdot a^v&=a^{u+v}, & \frac{a^u}{a^v}&=a^{u-v},\\
(a^u)^v&=a^{uv}, & a^{-u}&=\frac{1}{a^u},\\
\sqrt[n]{a^m}&=a^{m/n} &&(n\in\mathbb N,\ n\ge2).
\end{align*}
В задачах с отрицательным основанием и дробной степенью дополнительно учитывают область определения. Это особенно важно для уравнений с неизвестным основанием.

\chap{2. Показательные уравнения}
\key{Признак:} неизвестное находится в показателе степени, например $2^x=8$.

\subchap{Главное правило}
Если $a>0$ и $a\ne1$, то
\mymath{a^{f(x)}=a^{g(x)}\quad\Longleftrightarrow\quad f(x)=g(x).}
Поэтому сначала приводим обе части уравнения к \emph{одному основанию}, затем приравниваем показатели.

\example{Пример 1. Дробь и корень как степени}
\begin{align*}
2^x=\frac{1}{\sqrt[5]{8}}
&=\frac{1}{(2^3)^{1/5}}\\
&=2^{-3/5}
\quad\Longrightarrow\quad \boxed{x=-\frac35}.
\end{align*}

\clearpage
\example{Пример 2. Лишний множитель перед степенью}
\begin{align*}
8\cdot2^x&=\frac{1}{\sqrt[5]{16}},\\
2^3\cdot2^x&=2^{-4/5},\\
2^{x+3}&=2^{-4/5},\\
x+3&=-\frac45, \qquad
\boxed{x=-\frac{19}{5}}.
\end{align*}
\key{Проверка:} после подстановки найденного $x$ показатели слева и справа должны совпасть.

\chap{3. Степенные уравнения}
\key{Признак:} неизвестное находится в основании степени, например $x^{3/5}=8$.

\subchap{Как убрать показатель}
Когда допустимо, обе части возводят в степень, обратную исходной: $\frac35\cdot\frac53=1$. Но для отрицательных $x$ и чётных показателей необходимо отдельно проверить \emph{все} корни.

\example{Пример 3. Нечётная степень корня}
\mymath{x^{3/5}=125\quad\Longrightarrow\quad
x=125^{5/3}=(5^3)^{5/3}=5^5=3125.}
\example{Пример 4. Отрицательный показатель}
\mymath{x^{-3/5}=8\quad\Longrightarrow\quad
x=8^{-5/3}=(2^3)^{-5/3}=\frac1{32}.}
Условие $x\ne0$ выполнено.

\clearpage
\chap{4. Как не потерять отрицательный корень}
\key{Внимание:} если $x^{4/3}=16$, преобразование $\bigl(x^{4/3}\bigr)^{3/4}=x$ неверно для $x<0$. На самом деле результатом будет $|x|$.
\mymath{x^{4/3}=16
\quad\Longleftrightarrow\quad
\bigl(\sqrt[3]{x}\bigr)^4=2^4
\quad\Longleftrightarrow\quad
\sqrt[3]{x}=\pm2.}
Следовательно,
\mymath{\boxed{x=-8\quad\text{или}\quad x=8}.}
Оба числа подходят: $\left(\sqrt[3]{-8}\right)^4=(-2)^4=16$.

Аналогичная ситуация возникает для $x^{-2}=4$:
\mymath{\frac1{x^2}=4\quad\Longleftrightarrow\quad
x^2=\frac14\quad\Longleftrightarrow\quad
\boxed{x=\pm\frac12}.}

\subchap{Короткая памятка об области определения}
\begin{itemize}
\item Для $x^{m/n}$ с несократимой дробью $m/n$ и \textbf{чётным} $n$ требуется $x\ge0$ (при отрицательной степени также $x\ne0$).
\item При \textbf{нечётном} $n$ отрицательные $x$ допустимы; если числитель $m$ чётный, могут возникать два противоположных корня.
\item При отрицательном показателе обязательно $x\ne0$.
\end{itemize}

\clearpage
\chap{5. Прикладная задача: работаем с формулой}
На занятии использовалась формула для силы натяжения:
\mymath{T=\rho gL^3-mg,}
где $T$~--- сила (Н), $\rho$~--- плотность (кг/м$^3$), $g$~--- ускорение (м/с$^2$), $m$~--- масса (кг), $L$~--- искомая длина (м). Для физических размеров $L\ge0$.

\example{Пример 5. Определить длину по заданной силе}
Пусть $\rho=1000$, $g=10$, $m=25$, $T=1000$. Тогда
\begin{align*}
1000&=1000\cdot10\cdot L^3-25\cdot10,\\
10000L^3&=1250,\\
L^3&=\frac{1250}{10000}=\frac18,\\
&\boxed{L=\sqrt[3]{\frac18}=0{,}5\ \text{м}}.
\end{align*}

\subchap{Что означает «не больше»}
Если требуется \textbf{наибольшая} длина при условии $T\le1000$, нужно сохранить смысл неравенства:
\mymath{10000L^3-250\le1000
\quad\Longleftrightarrow\quad L^3\le\frac18
\quad\Longleftrightarrow\quad 0\le L\le0{,}5.}
Искомое наибольшее значение $L_{\max}=0{,}5$ м достигается при равенстве $T=1000$. Подстановка предельного значения обоснована тем, что правая часть формулы возрастает с ростом $L\ge0$.

\subchap{Мини-проверка}
\begin{enumerate}
\item Определили, где находится неизвестное: в основании, в показателе или в формуле?
\item Верно преобразовали корни, дроби, составные числа и показатели?
\item Учли область определения и возможные два корня?
\item Сохранили смысл слов «не больше» / «не меньше»?
\item Проверили ответ подстановкой и единицы измерения?
\end{enumerate}

\clearpage
\thispagestyle{plain}
\begin{center}
{\LARGE\bfseries\color{accentdark}Алгоритм: уравнения со степенями}\par
\vspace{8mm}
\end{center}
\begin{center}
\begin{tikzpicture}[
  every node/.style={font=\small},
  beginend/.style={ellipse,draw=accent,line width=.9pt,minimum width=46mm,minimum height=11mm,align=center},
  process/.style={rectangle,rounded corners=1mm,draw=accent,line width=.9pt,text width=70mm,minimum height=17mm,align=center,inner sep=3mm},
  narrow/.style={rectangle,rounded corners=1mm,draw=accent,line width=.9pt,text width=51mm,minimum height=21mm,align=center,inner sep=2.5mm},
  decision/.style={diamond,aspect=2.2,draw=accent,line width=.9pt,text width=40mm,align=center,inner sep=1.8mm},
  go/.style={-{Stealth[length=2.7mm]},line width=.9pt,draw=accentdark}
]
\node[beginend] (start) at (0,0) {Уравнение со степенями};
\node[process] (scan) at (0,-1.9) {Определите, где находится неизвестное};
\node[decision] (branch) at (0,-4.0) {В показателе?};
\node[narrow] (exp1) at (-4.0,-6.3) {Приведите обе части к одному основанию $a>0$, $a\ne1$};
\node[narrow] (pow1) at (4.0,-6.3) {Определите область допустимых значений и знаки};
\node[narrow] (exp2) at (-4.0,-9.2) {Приравняйте показатели и решите уравнение};
\node[narrow] (pow2) at (4.0,-9.2) {Выразите $x$; проверьте, нужны ли два корня};
\node[process] (check) at (0,-12.0) {Проверьте все корни подстановкой в исходное уравнение};
\node[beginend] (finish) at (0,-14.0) {Запишите ответ};
\draw[go] (start)--(scan);
\draw[go] (scan)--(branch);
\draw[go] (branch.south west) |- node[pos=.28,above left]{да} (exp1.north);
\draw[go] (branch.south east) |- node[pos=.28,above right]{нет} (pow1.north);
\draw[go] (exp1)--(exp2);
\draw[go] (pow1)--(pow2);
\draw[go] (exp2.south) -- ++(0,-.65) -| (check.west);
\draw[go] (pow2.south) -- ++(0,-.65) -| (check.east);
\draw[go] (check)--(finish);
\end{tikzpicture}
\end{center}
\vspace{6mm}
\begin{center}\color{muted}\small В прикладной задаче сначала составьте уравнение или неравенство из данных,\par затем выразите искомую величину и проверьте физический смысл.\end{center}

\clearpage
\chap{6. Тренировка --- Путь А}
\textbf{10 заданий.} Выполняйте решения письменно: фиксируйте приведение к общему основанию, раскрытие степеней и ограничения. В задачах 9--10 считайте, что $L\ge0$.

\begin{enumerate}
\item Решите уравнение: \(3^x=81.\)
\item Решите уравнение: \(2^x=\dfrac{1}{\sqrt[5]{32}}.\)
\item Решите уравнение: \(16\cdot2^x=\dfrac{1}{\sqrt[3]{4}}.\)
\item Решите уравнение: \(9\cdot3^{2x-1}=\dfrac{1}{\sqrt[3]{27}}.\)
\item Найдите все действительные корни: \(x^{3/5}=8.\)
\item Найдите все действительные корни: \(x^{5/2}=32.\)
\item Найдите все действительные корни: \(x^{-3/5}=8.\)
\item Найдите все действительные корни: \(x^{4/3}=81.\)
\item По формуле \(T=\rho gL^3-mg\) определите $L$, если
$T=390$ Н, $\rho=1000$ кг/м$^3$, $g=10$ м/с$^2$, $m=25$ кг.
\item Сила натяжения вычисляется по формуле \(T=\rho gL^3-mg\).
При $\rho=1000$ кг/м$^3$, $g=10$ м/с$^2$ и $m=10$ кг она должна быть \emph{не более} $2060$ Н. Найдите наибольшее допустимое значение $L$ в метрах.
\end{enumerate}

\vspace{4mm}
\textbf{Вопросы для самопроверки:} почему в № 8 могут получиться два ответа? Почему в № 10 нужно сначала записать неравенство?

\clearpage
\chap{7. Ответы}
\renewcommand{\arraystretch}{1.5}
\begin{center}
\begin{tabularx}{.73\linewidth}{@{}>{\bfseries}c X@{}}
\toprule
№ & Ответ\\
\midrule
1 & $x=4$\\
2 & $x=-1$\\
3 & $x=-\dfrac{14}{3}$\\
4 & $x=-1$\\
5 & $x=32$\\
6 & $x=4$\\
7 & $x=\dfrac1{32}$\\
8 & $x=-27$; $x=27$\\
9 & $L=0{,}4$ м\\
10 & $L_{\max}=0{,}6$ м\\
\bottomrule
\end{tabularx}
\end{center}

\vspace{10mm}
\textbf{Проверка ответов на прикладные задачи:}
\begin{align*}
\text{№ 9:}\quad 1000\cdot10\cdot(0{,}4)^3-25\cdot10
&=640-250=390\ \text{Н};\\
\text{№ 10:}\quad 1000\cdot10\cdot(0{,}6)^3-10\cdot10
&=2160-100=2060\ \text{Н}.
\end{align*}

При несовпадении ответа сначала проверьте знаки показателей и арифметику дробей, затем ограничения и возможную потерю отрицательного корня. Решение с корректными преобразованиями важнее одного совпавшего числа.

\vfill
\begin{center}{\color{muted}\small Лёвин Артём Александрович эксклюзивно для Марка Гукина}\end{center}
\end{document}
